\chapter{Declaración de uso de herramientas de inteligencia artificial}
\label{ap:declaracion-ia}

En este apéndice se documenta, de forma transparente, el uso que se ha hecho de herramientas de inteligencia artificial generativa durante la realización de este Trabajo de Fin de Grado.

\section{Alcance del uso}
\label{sec:ia_alcance}

Las herramientas de inteligencia artificial se han empleado como \textbf{apoyo}, y nunca como sustituto del criterio técnico del autor ni del trabajo de análisis, diseño y validación. Entre otras tareas, su uso ha incluido:

\begin{itemize}
  \item \textbf{Generación inicial de código:} obtención de borradores o esqueletos de scripts y componentes que después se han revisado, adaptado e integrado manualmente.
  \item \textbf{Organización de la documentación:} ayuda para estructurar notas, índices y resúmenes del material de trabajo.
  \item \textbf{Mejora de la redacción técnica:} sugerencias de estilo y claridad sobre textos ya escritos por el autor, sin delegar en la herramienta la elaboración del contenido.
  \item \textbf{Interpretación preliminar de patrones:} formulación de hipótesis iniciales sobre el comportamiento del tráfico, que posteriormente se han comprobado y validado con código y datos reales.
\end{itemize}

\section{Supervisión y validación}
\label{sec:ia_supervision}

Todo el contenido generado o sugerido con apoyo de estas herramientas ha sido \textbf{revisado, adaptado, ejecutado y validado por el autor}. En particular:

\begin{itemize}
  \item Las decisiones técnicas, el diseño del clasificador y la interpretación final de los resultados son responsabilidad del autor.
  \item El código se ha probado y comprobado sobre los datos del proyecto antes de darlo por válido.
  \item Las hipótesis sugeridas por las herramientas no se han aceptado de forma automática: se han contrastado de manera programática y solo se han incorporado cuando los datos las respaldaban.
  \item La detección de anomalías la realiza siempre el clasificador basado en reglas conductuales. Las herramientas de inteligencia artificial no toman decisiones de detección.
\end{itemize}

El apoyo de estas herramientas fue más amplio que esa ayuda inicial. También sirvieron para generar y adaptar los scripts de comparación con los modelos clásicos de aprendizaje automático, configurar el procesamiento de los datos y preparar el código auxiliar necesario para ejecutar los experimentos. Y ayudaron a revisar errores y resultados, a apoyar la implementación de la aplicación web y a mejorar la documentación, las figuras y la redacción de la memoria.

Ese apoyo se mantuvo siempre bajo el control del autor. Las herramientas ayudaron a producir y a revisar código, pero los scripts se ejecutaron en el entorno real del proyecto y sobre los datos reales del trabajo. El autor comprobó las salidas, contrastó las métricas y decidió en cada caso qué resultados se incorporaban a la memoria. Las herramientas de inteligencia artificial no realizaron por sí solas la evaluación académica ni sustituyeron la supervisión del autor y de los tutores.

En ningún caso debe entenderse que el trabajo se ha realizado de forma autónoma por una herramienta de inteligencia artificial. Su papel ha sido el de un \textbf{apoyo supervisado}, y la autoría, el criterio y la validación del Trabajo de Fin de Grado corresponden al autor.
